\chapter{Rendering, Auslieferung, Build und Betrieb}\label{kap:betrieb}

\section{Rendering}

Der Seitengenerator (\datei{RCode/generate_static_site.R}) baut die Seiten
mit dem R-Paket \code{htmltools}. Der Ablauf:

\begin{enumerate}
  \item Nur Ligen, für die alle benötigten Ergebnisse vorliegen, werden
        gerendert. Die übrigen werden übersprungen und im Log genannt: Eine
        Website ohne eine Liga ist besser als gar keine.
  \item Zuerst werden die Verlaufsdaten aller Ligen berechnet. Davon hängen
        der Link auf der Ligaseite und die Ziele der Navigation auf den
        Verlaufsseiten ab.
  \item Danach entstehen die Ligaseiten, die Verlaufsseiten (jede einzeln
        abgesichert), die Aufstiegsseite (nur wenn mindestens eine
        Regionalliga dabei ist) und die Methodik-Seite.
  \item Abschließend steht im Log \enquote{wrote 22 pages}.
\end{enumerate}

\paragraph{Atomares Schreiben.} Jede Seite wird zuerst in eine Datei
\code{<name>.tmp} im selben Verzeichnis geschrieben und dann umbenannt
(\datei{generate_static_site.R:960--980}). Das Umbenennen ist auf demselben
Dateisystem atomar. Der Webserver sieht deshalb immer entweder die alte oder
die vollständige neue Seite, nie eine halb geschriebene.

\paragraph{Assets.} Stilvorlage, Schriften, Favicon und \code{verlauf.js}
werden bei jedem Rendern nach \code{assets/} kopiert.

\section{Auslieferung}

Caddy bindet das Volume nur lesend ein, komprimiert mit zstd und gzip und
setzt Cache-Regeln je Pfad:

\begin{center}
\small
\begin{tabular}{@{}ll@{}}
\toprule
Pfad & \code{Cache-Control} \\
\midrule
Ligaseiten, Verlaufsseiten, Aufstiegsseite & \code{no-cache} \\
\code{methodik.html} & \code{public, max-age=3600} \\
\code{/assets/*} & \code{public, max-age=300} \\
\bottomrule
\end{tabular}
\end{center}

\begin{achtung}[Neue Seiten brauchen eine eigene Caddy-Zeile]
Die Pfadregeln in Caddy gelten exakt. Eine Seite, die im \code{route}-Block
nicht genannt ist, bekommt \emph{gar keinen} Cache-Header. Wer eine Seite
hinzufügt, trägt sie vor dem Deployment in die Caddy-Konfiguration ein. So
geschah es zuletzt am 11.~Oktober 2026 für die zehn Verlaufsseiten.
\end{achtung}

\section{Container}

Das Produktions-Image (\datei{Dockerfile}) wird in drei Stufen gebaut:

\begin{enumerate}
  \item \code{rust:1.98-alpine} kompiliert den Rust-Server.
  \item \code{rocker/r-ver:4.6.1} installiert die R-Pakete mit allen
        Entwicklungsbibliotheken.
  \item Das Laufzeit-Image \code{rocker/r-ver:4.6.1} enthält nur die nötigen
        Bibliotheken, die R-Pakete, das Rust-Programm und \datei{RCode/}.
        Es läuft als Benutzer \code{appuser} (uid 1001), nicht als root.
\end{enumerate}

Die R-Pakete im Produktionspfad sind \code{dplyr}, \code{tidyr},
\code{tibble}, \code{magrittr}, \code{httr}, \code{jsonlite}, \code{digest},
\code{xtable}, \code{here} und \code{htmltools}. Tests und Skripte gehören
nicht zum Image.

\begin{center}
\small
\begin{tabularx}{\textwidth}{@{}lX@{}}
\toprule
Einstellung & Wert \\
\midrule
Healthcheck & \code{curl -f http://localhost:8080/health}, alle 30~s \\
Umgebung & \code{RAPIDAPI\_KEY} (Pflicht), \code{SEASON=2026}, \code{DURATION} leer, \code{TZ=Europe/Berlin}, \code{STATIC\_SITE\_DIR=/app/ShinyApp/public} \\
Volume & \code{fussball-site} $\rightarrow$ \code{/app/ShinyApp/public} \\
Ressourcen & 2~GB Speicher, 4 CPUs, Logs $3 \times 10$~MB \\
Neustart & \code{unless-stopped} \\
\bottomrule
\end{tabularx}
\end{center}

Der Healthcheck prüft nur, ob der Rust-Server lebt. Ein hängender
R-Scheduler macht den Container nicht \enquote{unhealthy}. Nach einem
Deployment muss man deshalb die Logs lesen, nicht nur den Status.

\section{Continuous Integration}

GitHub Actions (\datei{.github/workflows/ci.yml}) prüft jeden Pull Request
und jeden Push auf \code{main}:

\begin{description}[font=\sffamily\bfseries\small]
  \item[rust-quality] \code{cargo fmt --check}, \code{cargo clippy -D warnings},
        \code{cargo test --release} (76 Tests).
  \item[r-lint] \code{lintr} 3.4.0 (fest gepinnt) über \datei{RCode/}, mit
        Zeilenlänge 120. Jeder Befund lässt den Job scheitern.
  \item[image-build-and-test] baut das Produktions-Image, startet darin den
        Rust-Server und führt den vollständigen R-Testsatz aus, einschließlich
        der JavaScript-Tests mit Node 24 und jsdom. Ein Test, der übersprungen
        wird, weil Rust oder Node fehlt, gilt in der CI als Fehler.
  \item[push-image] nur auf \code{main}: Das Image geht mit den Tags
        \code{:latest} und \code{:<kurzer Commit-SHA>} zu Docker Hub
        (\code{chrisschwer/league-simulator}).
  \item[report-failure] legt bei einem Fehlschlag auf \code{main} ein Issue an.
\end{description}

Dazu kommt CodeQL, das wöchentlich und bei jedem Pull Request die
Workflow-Dateien prüft.

\section{Tests}

\begin{itemize}
  \item \textbf{R}: rund 870 \code{testthat}-Blöcke in 61 Dateien. Die
        Dateien folgen einer festen Namensregel
        \code{test-<einheit>[-<thema>].R}, wobei \code{<einheit>} der Name
        einer Datei in \datei{RCode/} ist. Ein Wächtertest erzwingt die Regel.
  \item \textbf{JavaScript}: Sortierung, Tooltips, Veraltet-Hinweis und das
        Verlaufsdiagramm werden in Node mit jsdom getestet, aufgerufen aus den
        R-Tests der erzeugenden Einheit.
  \item \textbf{Rust}: 76 Tests, darunter die Erhaltung der ELO-Summe und
        die Gleichheit von Tormodell in Simulation und Spieldetails.
  \item \textbf{Naht R–Rust}: Integrationstests gegen einen laufenden
        Rust-Server prüfen die JSON-Verträge.
\end{itemize}

\section{Betrieb}

Produktiv läuft der Container auf einem eigenen Server. Dessen
\code{compose.yml} pinnt das Image auf den Commit-SHA eines Merges auf
\code{main}, nie auf \code{latest}, und es gibt keinen automatischen
Updater. Ein Deployment ist im Kern eine einzige geänderte Zeile:

\begin{enumerate}
  \item Umfang prüfen: was sich seit dem laufenden Stand geändert hat, laut
        \code{git log} und nicht laut Pull-Request-Text.
  \item Neue Seiten vorab in die Caddy-Konfiguration eintragen.
  \item Pin auf den neuen SHA setzen, \code{docker compose pull} und
        \code{up -d}.
  \item Verifizieren, und zwar bis über die zweite Runde des Schedulers
        hinaus: kein Neustart-Kreislauf, \enquote{wrote 22 pages}, keine
        Warnungen im Log, alle Seiten mit Status 200 und korrektem
        Cache-Header.
\end{enumerate}

Ein Rollback setzt den Pin auf den vorherigen SHA zurück. Das alte Image
liegt lokal auf dem Server.

\section{Werkzeuge für Betreiber}

\begin{center}
\small
\begin{tabularx}{\textwidth}{@{}lX@{}}
\toprule
Skript & Zweck \\
\midrule
\code{scripts/season\_transition.R} & Saisonwechsel: End-ELO berechnen, neue Stammdatei erzeugen \\
\code{scripts/preview\_site.R} & lokale Vorschau der Website aus einer gespeicherten Ergebnisdatei \\
\code{scripts/verlauf\_fixture.R} & Beispieldaten für die Vorschau der Verlaufsseiten (braucht den Rust-Server) \\
\code{scripts/calibrate\_historical\_elo.R} & einmalige ELO-Eichung neuer Ligen, schreibt nur mit \code{--confirm} \\
\code{scripts/analyze\_league\_empirics.R} & Empirie-Report: Tore je Spiel, Heim/Remis/Auswärts, Remisquoten \\
\bottomrule
\end{tabularx}
\end{center}

Destruktive Werkzeuge laufen standardmäßig als Probelauf und verändern erst
mit \code{--confirm} etwas.
